% Part: sets-functions-relations
% Chapter: sets
% Section: basics

\documentclass[../../../include/open-logic-section]{subfiles}

\begin{document}

\olfileid{sfr}{set}{bas}
\olsection{عنصراں نال تعیّن}

\emph{سیٹ} شےواں دا اک اکٹھ اے، جہنوں {اکو} شے منّیا جاندا اے۔ سیٹ بناؤن والیاں
شےواں نوں اوس سیٹ دے \emph{عنصر} یا \emph{رکن} آکھیا جاندا اے۔ جے $x$ سیٹ~$A$ دا
!!a{element} ہووے، تاں اسی $x \in A$ لکھدے آں؛ جے نہ ہووے، تاں $x \notin A$
لکھدے آں۔ جس سیٹ دے کوئی !!{element}s نہ ہون، اوہنوں \emph{خالی} سیٹ آکھدے نیں
تے اوہنوں~``$\emptyset$'' نال ظاہر کردے نیں۔

\begin{explain}
ایس گل نال کوئی فرق نہیں پیندا کہ اسی سیٹ نوں کِویں \emph{دسدے} آں، اوہدے
\pnboblique{!!{element}s} نوں کِہڑی \emph{ترتیب} وچ رکھدے آں، یا اوہدے \pnboblique{!!{element}s} نوں
\emph{کِنّی واری} گِندے آں۔ اصل گل صرف ایہ اے کہ اوہدے !!{element}s کِہڑے نیں۔
اسی ایس گل نوں ہِیٹھلے اصول دی صورت وچ بیان کردے آں۔
\end{explain}

\begin{defn}[عنصراں نال تعیّن]
  جے $A$ تے $B$ سیٹ نیں، تاں $A = B$ اُدوں تے صرف اُدوں ہوندا اے جدوں
  $A$ دا ہر !!{element}، $B$ دا وی !!a{element} ہووے، تے ایس توں اُلٹ وی۔
\end{defn}

عنصراں نال تعیّن دا اصول کجھ علامتاں ورتن دا جواز دیندا اے۔ عام طور تے، جدوں
ساڈے کول کجھ شےواں $a_{1}$، \dots، $a_{n}$ ہون، تاں $\{a_{1}, \dots, a_{n}\}$
\emph{اوہی} سیٹ اے جہدے !!{element}s، $a_1, \ldots, a_n$ نیں۔ اسی لفظ
``\emph{اوہی}'' اُتے زور دیندے آں، کیوں جو عنصراں نال تعیّن دا اصول دسدا اے کہ
ایہو جہا سیٹ صرف \emph{اک} ہو سکدا اے۔ ایس اصول توں ایہ وی جائز اے:
  \[
    \{a, a, b\} = \{a, b\} = \{b,a\}.
  \]
ایس نال اوہ گل واضح ہوندی اے کہ سیٹاں بارے سوچدیاں سانوں ایہدی پرواہ نہیں
ہوندی کہ اوہناں دے !!{element}s دی ترتیب کی اے، یا اوہ کِنّی واری دسے گئے نیں۔

\begin{tagblock}{novice}
\begin{ex}
جدوں وی تہاڈے کول کجھ شےواں ہون، تسیں اوہناں نوں اک سیٹ وچ اکٹھیاں کر سکدے او۔
مثال لئی، رچرڈ دے بھین بھراواں دا سیٹ اک ایہو جہا سیٹ اے جہدے وچ اک بندہ اے،
تے اسی اوہنوں $S=\{\textrm{Ruth}\}$ لکھ سکدے آں۔ $4$ توں گھٹ مثبت صحیح
اعداد دا سیٹ $\{1, 2, 3\}$ اے، پر اوہنوں $\{3, 2, 1\}$ یا
$\{1, 2, 1, 2, 3\}$ وی لکھیا جا سکدا اے۔ عنصراں نال تعیّن دے اصول موجب ایہ
سارے اکو سیٹ نیں۔ کیوں جو $\{1, 2, 3\}$ دا ہر !!{element}،
$\{3, 2, 1\}$ دا (تے $\{1, 2, 1, 2, 3\}$ دا وی) !!a{element} اے، تے ایس
توں اُلٹ وی۔
\end{ex}
\end{tagblock}

اکثر اسی سیٹ نوں اوہدے !!{element}s دی سانجھی خاصیت نال دسّاں گے۔ ایس لئی
اسی ایہ مختصر علامتی لکھت ورتاں گے: $\Setabs{x}{\phi(x)}$، جتھے $\phi(x)$
اوہ خاصیت اے جہڑی $x$ وچ ہونی چاہیدی اے تاں جو اوہ سیٹ دے \pnboblique{!!{element}s} وچوں
اک گنیا جا سکے۔

\begin{tagblock}{novice}
\begin{ex}
اپنی مثال وچ اسی $S$ نوں ایویں وی دس سکدے ساں:
\[
S = \Setabs{x}{x \text{ رچرڈ دی بھین یا اوہدا بھرا اے}}.
\]
\end{ex}
\end{tagblock}

\begin{tagblock}{math}
\begin{ex}
کِسے عدد نوں \emph{کامل} اُدوں تے صرف اُدوں آکھدے نیں جدوں اوہ اپنے اوہناں
ونڈن والے اعداد دے جوڑ دے برابر ہووے جو عدد نوں پورا ونڈدے ہون پر آپ اوہدے
برابر نہ ہون۔ مثال لئی، $6$ کامل اے، کیوں جو اوہنوں ونڈن والے اعداد، جو آپ
اوہدے برابر نہیں، $1$، $2$ تے~$3$ نیں، تے $6 = 1 + 2 + 3$ اے۔ دراصل،
$6$، $10$ توں گھٹ اکو اک مثبت صحیح عدد اے جو کامل اے۔ سو، عنصراں نال تعیّن
دا اصول ورت کے اسی آکھ سکدے آں:
\[
	\{6\} = \Setabs{x}{x\text{ کامل اے تے }0 \leq x \leq 10}
\]
سجے پاسے دی علامتی لکھت نوں اسی ایویں پڑھدے آں: ``اوہناں $x$ دا سیٹ جیہناں لئی
$x$ کامل اے تے $0 \leq x \leq 10$''۔ ایتھے برابری دی ایہ لکھت تصدیق کردی
اے کہ سیٹاں بارے سوچدیاں سانوں ایہدی پرواہ نہیں ہوندی کہ اوہناں نوں کِویں
دسیا گیا اے۔ ہور عام طور تے، عنصراں نال تعیّن دا اصول ایہ یقینی بناندا اے کہ
اوہناں $x$ دا سیٹ جنہاں لئی $\phi(x)$ ہووے، ہمیشہ اکو ہوندا اے۔
سو، عنصراں نال تعیّن دا اصول $\Setabs{x}{\phi(x)}$ نوں اوہناں $x$ دا
\emph{اوہی} سیٹ آکھن دا جواز دیندا اے جنہاں لئی~$\phi(x)$ ہووے۔
\end{ex}
\end{tagblock}

عنصراں نال تعیّن دا اصول ایہ وکھاؤن دا طریقہ دیندا اے کہ سیٹ اکو نیں: $A = B$
وکھاؤن لئی ایہ وکھاؤ کہ جدوں وی $x \in A$ ہووے، تاں $x \in B$ وی ہووے،
تے جدوں وی $y \in B$ ہووے، تاں $y \in A$ وی ہووے۔

\begin{prob}
ثابت کرو کہ ودھ توں ودھ اک خالی سیٹ ہو سکدا اے، یعنی ایہ وکھاؤ کہ جے $A$ تے
$B$ ایہو جہے سیٹ نیں جنہاں دے کوئی !!{element}s نہیں، تاں $A = B$ اے۔
\end{prob}

\end{document}
